\scanpage{12}
% Pencil foliation: 10, upper right; 11, lower left.
Noter que si $F$ est compatible avec $T^\ell_{\unsure{\xi|\zeta}}$, il est
à fortiori \uline{associé à} $\unsure{\xi|\zeta}$, et on trouve
\diagram{\begin{tikzcd}
G\deleted{\unsure{\mathbb{Q}}}\arrow[r]&\pi_1(S,\unsure{\xi|\zeta})
\end{tikzcd}}
d'où
\diagram{\begin{tikzcd}
\unsure{\psi_\ell|\eta_\ell}:G(\mathbb{Q}_\ell)\arrow[r]&\pi_1(S,\unsure{\xi|\zeta})
\end{tikzcd}}
On voit tout de suite que cet homomorphisme
est inverse à gauche de $\varphi_\ell$
\[
 \boxed{\unsure{\psi_\ell|\eta_\ell}\varphi_\ell
   =\operatorname{id}_{\pi_1(S,\unsure{\xi|\zeta})}}
\]
Si on s'intéresse à une sous-catégorie $C_\alpha$ de $C$,
correspondant à un $\unsure{G_\alpha^0|G_\alpha}$, alors $G_\alpha/G_\alpha^0$
\deleted{\unsure{correspond}}\inserted{définit}
\deleted{\unsure{quotient}} un quotient $\pi_\alpha$ de $\pi_1(S,\unsure{\xi|\zeta})$,
et on a
\diagram{\begin{tikzcd}
& G_\alpha/G_\alpha^0 \arrow[dr,dashed] &\\
\pi_1(S,\unsure{\xi|\zeta})\arrow[r,"\varphi_\ell"]
 \arrow[rr,bend right=32,swap,"\uline{\text{hom can.}}"]
 & G_\alpha(\mathbb{Q}_\ell)\arrow[ur,dashed]\arrow[r]&\pi_\alpha
\end{tikzcd}}
(commutatif).

La conjecture de Tate \inserted{(avec action semi-simple de Galois !)}
équivaut à dire que, si
$S$ est localisé d'un schéma de type fini sur $\mathbb{Z}$, alors
$G_\alpha\unsure{\otimes|\times}_{\mathbb{Q}_\ell}$ est
l'enveloppe algébrique de $\operatorname{Im}(\pi_1(\uline{S,\unsure{\xi|\zeta}})
\longrightarrow G_\alpha(\mathbb{Q}_\ell))$.

\margin{\uline{NB} La conjecture
de Tate sous la
forme : les \unsure{$\frac12$-invariants}
dans $H^{2i}(\overline X,\mathbb{Q}_\ell(i))$
sont les comb.
à coeff. dans $\mathbb{Q}_\ell$
des classes algébriques
--- implique (avec
conjectures standard)
la $\frac12$ simplicité dans
Tate ---}

\margin{\deleted{\fbox{\begin{tabular}{l}
\unsure{l'image} \illegible{4}\\
\unsure{$\pi_1$} \illegible{2} \unsure{$(M)$}\\
\illegible{4}
\end{tabular}}}}

11) \uline{Cas} $S=\operatorname{Spec}k$, $k$ \uline{corps fini}.

$C\subset \mathcal{M}(k)$, et $F$ un ``foncteur'' fibre sur
$C$ à coeff dans $K$. Comme on a un élément
canonique de
\diagram{\begin{tikzcd}
\operatorname{Frob}_k\in\operatorname{Aut}\underline{\operatorname{Id}}_{\mathcal{M}(k)}
 \arrow[r]&\operatorname{Aut}\operatorname{id}_C
 \left[\simeq\unsure{Z|\mathfrak Z}(\mathbb{Q})\right]
\end{tikzcd}}
\margin{\inserted{\unsure{$Z$} est le centre de $G$}\deleted{\unsure{$Z(K)$}}}
on en déduit un élément canonique de
$G(K)=\operatorname{Aut}F$, et même de $\unsure{Z|\mathfrak Z}(\mathbb{Q})$
($\unsure{Z|\mathfrak Z}$ est le centre de $G$)
\[
 \boxed{\operatorname{Frob}_{k,F}\in
 \unsure{Z|\mathfrak Z}(\mathbb{Q})\subset G(K)}.
\]
\margin{\fbox{\begin{tabular}{l}
NB. $\unsure{Z|\mathfrak Z}$ est\\
en fait défini\\
sur $\mathbb{Q}$ !
\end{tabular}}}
